\section{Panel: The AGenSE Alliance}
\label{sec:panel}

The workshop opened with the panel ``The AGenSE Alliance: Shaping the Future of Agentic Software Engineering'', moderated by Giuseppe Di Fatta. Horizon Europe projects presented their views on agentic software engineering: Ali Atalay (AI4SEC, coordinator of AI4SWEng), Anh Tuan Le (University of Warwick, ASTIR), Ioannis Kyprakis (FORTH, FAITH), and Jean-Marie Mottu (Nantes Université) and Tuong Manh Vu (University of York) for MOSAICO. \todo{HIVEMIND: represented or not?}

\textbf{AI4SWEng} is building an AI engineering suite that connects requirements, architecture, coding, assurance and deployment through a shared workflow with a human review gate, validated in five use cases; its evaluation tracks developer time, code quality, energy use, adoption and technical debt. \textbf{ASTIR} targets robotics, where agent mistakes move physical machines: specialised agents hand structured, traceable artefacts from requirements to runtime monitoring, deterministic checks run before an LLM is consulted, and engineers remain accountable for every decision. \textbf{FAITH} contributed a trustworthiness perspective through its AI Trustworthiness Assessment Framework, validated in seven large-scale pilots and being extended to LLM- and agent-specific threats such as prompt injection, excessive agency and over-reliance. \textbf{MOSAICO} presented communities of agents with explicit roles, workflows modelled in SysML~v2 and agent-to-agent protocols, weighing the benefits of specialisation and redundancy against coordination overhead and error propagation.

The discussion was organised around questions that recurred in the papers. How can agents follow intended changes to requirements and domain knowledge without drifting from the original intent, from requirements to validated code? How do we know that an agent's improvement is real, what should be measured before agents touch a system, and how do we exclude gains that only come from changed tests? What should a benchmark covering the full lifecycle look like, and should agentic workflows be compared with human teams, single agents or both? How much autonomy should agents have, who is accountable when an agent-built system fails in a safety-critical setting, and how can provenance be kept across hand-offs so that one agent's error is caught rather than passed on? What does trustworthiness mean beyond accuracy, in the light of regulation such as the EU AI Act? And when does dividing work among specialised agents beat one strong agent with good tools? The panel closed by asking which shared standard, benchmark or dataset the projects could offer the community. \todo{Main positions and points of agreement or disagreement from the discussion.}
